\chapter{Tilting Ball Maze}

\section{Overview}
In this project, you will build a handheld, tilt-controlled arcade game: the \textbf{Tilting Ball Maze}!
By wiring an MPU-6050 6-axis accelerometer and a 128$\times$64 OLED display to your ESP32-C3 over a shared
I\textsuperscript{2}C bus, your microcontroller will measure how much you tilt the breadboard in real time
and roll a ball sprite through a maze on screen. Guide the ball from the starting corner to the goal without
getting trapped in dead ends!

Over the course of this project, you will:
\begin{itemize}
    \item Connect multiple digital devices to a single shared I\textsuperscript{2}C bus
    \item Read tilt and gravity data from an MPU-6050 accelerometer
    \item Drive a high-resolution monochrome OLED screen using MicroPython's \texttt{framebuf} graphics library
    \item Implement 2D physics (acceleration, velocity, friction, and wall collisions) in Python
    \item Build a complete, responsive interactive game
\end{itemize}

\begin{figure}[H]
\centering
    \maybeincludegraphics[width=.6\linewidth]{project_6/success.jpg}
    \caption{The end result should look something like this}
\end{figure}

\begin{tcolorbox}[colback=blue!5!white,colframe=blue!40!black]
    \textbf{Learning goals:}
    \begin{itemize}
        \item Understand how the I\textsuperscript{2}C protocol uses two shared wires (SCL and SDA) to talk to multiple chips
        \item Measure physical tilt using 16-bit accelerometer readings
        \item Draw shapes, text, and sprites on an SSD1306 OLED display
        \item Calculate axis-aligned bounding box (AABB) collisions between a moving ball and maze walls
    \end{itemize}
\end{tcolorbox}

\pagebreak

\section{Components}

\begin{components}
    \item Seeed XIAO ESP32-C3 microcontroller
    \item 0.96-inch 128$\times$64 SSD1306 I\textsuperscript{2}C OLED display
    \item MPU-6050 6-axis accelerometer/gyroscope module (GY-521)
    \item Breadboard
    \item Jumper wires
\end{components}

\section{Understanding \texorpdfstring{I\textsuperscript{2}C}{I2C} Communication} \label{i2c_basics}

Up to now, every peripheral had its own dedicated GPIO pin. However, microcontrollers have a limited
number of pins. The \textbf{I\textsuperscript{2}C} (Inter-Integrated Circuit) communication protocol solves
this problem by allowing multiple sensor and display chips to share the exact same two wires:

\begin{itemize}
    \item \textbf{SCL (Serial Clock):} A clock signal that synchronizes data transmission between chips.
    \item \textbf{SDA (Serial Data):} A bidirectional data line that carries the messages.
\end{itemize}

Every I\textsuperscript{2}C device has a unique digital \textbf{address} (like a street address). When the
microcontroller sends data on the SDA wire, it prefixes the message with the target device's address.
All other devices on the bus ignore the message.

In this project, both the OLED display (address \texttt{0x3C}) and the MPU-6050 accelerometer (address
\texttt{0x68}) connect to the same pair of microcontroller pins: \textbf{GPIO 7} (SCL) and \textbf{GPIO 6} (SDA).

\begin{tcolorbox}[colback=yellow!10!white,colframe=yellow!50!black]
    \textbf{Important:} Power both the OLED and the MPU-6050 from \textbf{3.3V}, not 5V.
    The ESP32-C3 GPIO pins operate at 3.3V logic levels. Connecting 5V power to the sensors could damage
    the microcontroller.
\end{tcolorbox}

\section{Directions}

\subsubsection{Remove previous components}
Before beginning, remove any components and wires from previous projects. Leave the microcontroller
firmly seated in its standard position.

\subsection{Creating the circuit}

% Seating the microcontroller. This is identical for every project, so the
% coordinates live here rather than in any one project chapter. The diagram
% generator reads the \bbhole definitions below for every breadboard diagram
% it draws, which is why this file is the only place they appear.
\subsubsection{Attach the microcontroller to the breadboard}
Orient the microcontroller so that the pin labeled \textbf{5V} goes into hole \bbhole{mcu.5v}{H1}
(Column H, Row 1), \textbf{GPIO2} into \bbhole{mcu.gpio2}{D1}, \textbf{GPIO20} into
\bbhole{mcu.gpio20}{H7}, and \textbf{GPIO21} into \bbhole{mcu.gpio21}{D7}. Then press it down---this
takes more force than you might expect.

\begin{figure}[H]
    \centering
    \maybeincludegraphics[width=.95\linewidth]{common/microcontroller_seated_in_breadboard.png}
    \caption{So far, so good!}
\end{figure}


\subsubsection{MPU-6050 accelerometer placement and wiring}
The MPU-6050 (GY-521) breakout board connects with straight-through header pins and sits flat facing upward.
Push the module's eight pins into column \textbf{A} (rows 12--19) on the front of the breadboard:
\textbf{INT} into \bbhole{mpu.header.int}{A12},
\textbf{AD0} into \bbhole{mpu.header.ad0}{A13},
\textbf{XCL} into \bbhole{mpu.header.xcl}{A14},
\textbf{XDA} into \bbhole{mpu.header.xda}{A15},
\textbf{SDA} into \bbhole{mpu.header.sda}{A16},
\textbf{SCL} into \bbhole{mpu.header.scl}{A17},
\textbf{GND} into \bbhole{mpu.header.gnd}{A18}, and
\textbf{VCC} into \bbhole{mpu.header.vcc}{A19}.

Only four pins (VCC, GND, SCL, SDA) are needed. Pins INT, AD0, XCL, and XDA remain unconnected.
Connect the MPU-6050 jumpers into the corresponding holes in row \textbf{E}:

\begin{center}
\begin{tabular}{|l|l|l|}
    \hline
    \textbf{MPU-6050 pin} & \textbf{XIAO pin} & \textbf{Jumper holes} \\
    \hline
    \connectionrow{mpu-power}{VCC}{3V3}{J3}{E19}
    \connectionrow{mpu-ground}{GND}{GND}{J2}{E18}
    \connectionrow{mpu-scl}{SCL}{GPIO 7}{B6}{E17}
    \connectionrow{mpu-sda}{SDA}{GPIO 6}{B5}{E16}
    \hline
\end{tabular}
\end{center}

\subsubsection{OLED display placement and wiring}
The 0.96-inch SSD1306 OLED module also connects with straight-through header pins, sitting flat facing upward.
Insert the four pins into column \textbf{A} (rows 25--28) next to the accelerometer:
\textbf{GND} into \bbhole{oled.header.ground}{A25},
\textbf{VCC} into \bbhole{oled.header.power}{A26},
\textbf{SCL} into \bbhole{oled.header.scl}{A27}, and
\textbf{SDA} into \bbhole{oled.header.sda}{A28}.

\begin{tcolorbox}[colback=yellow!10!white,colframe=yellow!50!black]
    \textbf{Check your OLED pins:} Most OLED displays have the pin order \textbf{GND, VCC, SCL, SDA}.
    However, some modules swap \textbf{VCC} and \textbf{GND}. Always check the silkscreen labels printed on
    your specific OLED board before powering on!
\end{tcolorbox}

Connect the OLED by daisy-chaining power, ground, clock, and data directly from the accelerometer:

\begin{center}
\begin{tabular}{|l|l|l|}
    \hline
    \textbf{OLED pin} & \textbf{Connected to} & \textbf{Jumper holes} \\
    \hline
    \connectionrow{oled-ground}{GND}{MPU GND}{D18}{E25}
    \connectionrow{oled-power}{VCC}{MPU VCC}{D19}{E26}
    \connectionrow{oled-scl}{SCL}{MPU SCL}{D17}{E27}
    \connectionrow{oled-sda}{SDA}{MPU SDA}{D16}{E28}
    \hline
\end{tabular}
\end{center}

All four connections are daisy-chained through the accelerometer rows. Its \textbf{VCC} and \textbf{GND}
connections continue from rows 19 and 18 to the OLED at rows 26 and 25. Likewise, \textbf{GPIO 7}
connects to the accelerometer's SCL pin at row 17 and continues to the OLED's SCL pin at row 27, while
\textbf{GPIO 6} connects to the accelerometer's SDA pin at row 16 and continues to the OLED's SDA pin at
row 28. Both devices therefore share the same power and I\textsuperscript{2}C bus connections.

\begin{figure}[H]
    \centering
    \maybeincludegraphics[width=.95\linewidth]{project_6/wiring.png}
    \caption{Connect both the OLED display and MPU-6050 accelerometer to the shared I2C bus.}
\end{figure}

\subsection{Programming the microcontroller}

Upload the following files to your microcontroller using the IDE:
\begin{itemize}
    \item \texttt{project\_6\_tilting\_ball\_maze.py}
    \item \texttt{lib/ssd1306.py} (in the \texttt{lib/} directory)
    \item \texttt{lib/mpu6050.py} (in the \texttt{lib/} directory)
\end{itemize}

Once the files are uploaded, open \texttt{project\_6\_tilting\_ball\_maze.py} and click the play button.
Refer back to Chapter \ref{ide} for IDE instructions.

\subsection{How the program works}

\subsubsection{Initializing the Shared I\textsuperscript{2}C Bus}
The program begins by creating a hardware I\textsuperscript{2}C bus on GPIO 7 (SCL) and GPIO 6 (SDA), then
scans for connected devices:

\begin{lstlisting}[language=Python, caption=Initializing I2C and discovering devices]
i2c = I2C(0, scl=Pin(7), sda=Pin(6), freq=400000)
devices = i2c.scan()
print("I2C devices found:", [hex(d) for d in devices])
\end{lstlisting}

If both modules are wired correctly, the serial console prints \texttt{['0x3c', '0x68']}.

\subsubsection{Reading Tilt from the Accelerometer}
The MPU-6050 contains micro-electromechanical sensors (MEMS) that measure acceleration in three dimensions ($X$, $Y$, and $Z$).
When the breadboard is held flat, gravity points straight down along the $Z$-axis. When you tilt the board,
gravity tilts into the $X$ and $Y$ axes:

\begin{lstlisting}[language=Python, caption=Converting tilt into ball acceleration]
ax, ay, az = mpu.read_accel()

# Remove the neutral offsets measured while the board was level
ax -= ax_offset
ay -= ay_offset

# Apply a small deadzone so the ball rests still on flat surfaces
if abs(ax) < DEADZONE:
    ax = 0
if abs(ay) < DEADZONE:
    ay = 0

# Integrate acceleration into velocity with friction
vx = (vx + ax * TILT_SCALE) * FRICTION
vy = (vy + ay * TILT_SCALE) * FRICTION
\end{lstlisting}

Accelerometers usually report a small nonzero value even when perfectly level. When the program starts,
it displays \textbf{Hold level...} and averages 100 readings to find that neutral X/Y offset. Keep the
breadboard level and still until the game title appears. The program subtracts those offsets from every
later reading so motion should be equally responsive in opposite directions.

\subsubsection{Collision Detection and Drawing}
The maze consists of rectangular wall segments. Before updating the ball's position, the game tests whether
the new position would overlap any wall. If a collision is detected on an axis, the ball stops against the wall:

\begin{lstlisting}[language=Python, caption=Axis-aligned bounding box collision]
new_x = ball_x + vx
if not hits_wall(new_x, ball_y, BALL_RADIUS):
    ball_x = new_x
else:
    vx = 0  # Stop horizontal motion on impact
\end{lstlisting}

Every frame, the program clears the screen buffer, draws all maze walls, renders the checkered goal zone,
draws the ball sprite, and calls \texttt{oled.show()} to push the entire image to the OLED over I\textsuperscript{2}C.

\begin{tcolorbox}[colback=yellow!10!white,colframe=yellow!50!black]
    \textbf{If something does not work:}
    \begin{itemize}
        \item \textbf{I\textsuperscript{2}C scan is empty:} Check the power (3.3V) and ground jumpers. Verify that SCL is on
        GPIO 7 (\bbref{jumper.oled-scl.start} and \bbref{jumper.mpu-scl.start}) and SDA is on GPIO 6
        (\bbref{jumper.oled-sda.start} and \bbref{jumper.mpu-sda.start}).
        \item \textbf{OLED remains dark:} Check whether your OLED module has GND and VCC in reversed order.
        Verify that \texttt{lib/ssd1306.py} was uploaded to the \texttt{lib/} folder.
        \item \textbf{Ball moves backwards:} Tilting forwards should roll the ball down; tilting right should roll it right.
        If an axis feels inverted, change \texttt{INVERT\_X} or \texttt{INVERT\_Y} from \texttt{False} to \texttt{True} at the
        top of the script.
        \item \textbf{Ball responds too slowly:} Increase \texttt{TILT\_SCALE} slightly. If very small tilts are ignored,
        decrease \texttt{DEADZONE} slightly.
        \item \textbf{Ball drifts or responds differently in opposite directions:} Restart the program and keep the
        breadboard level and still during calibration. If it still drifts, slightly increase \texttt{DEADZONE}.
    \end{itemize}
\end{tcolorbox}

\section{Review}
In this project, you combined motion sensing and display hardware on a single bus. You learned:
\begin{itemize}
    \item How the I\textsuperscript{2}C protocol connects multiple sensors and displays to just two microcontroller pins
    \item How an accelerometer detects gravitational tilt
    \item How to use MicroPython framebuffer commands (\texttt{fill}, \texttt{rect}, \texttt{fill\_rect}, \texttt{text})
    \item How basic physics simulation (velocity, dampening friction, and bounding box collisions) creates fluid game movement
\end{itemize}

\section{Possible Extensions}
\begin{itemize}
    \item \textbf{Design custom mazes:} Modify the \texttt{WALLS} list in the Python script to create your own challenging labyrinths.
    \item \textbf{Add hazard holes:} Add circular ``pit'' zones that send the ball back to the starting point if touched!
    \item \textbf{Add a timer:} Use \texttt{utime.ticks\_ms()} to track how fast players solve the maze and display their best time.
    \item \textbf{Multi-level maze:} When the goal is reached, load a new, more difficult maze layout.
    \item \textbf{Add sound effects:} Reconnect the speaker from Projects 3 and 4 to play a click when bumping into walls and a fanfare on victory!
\end{itemize}
